%! AP Statistics — Unit 1, Lesson 1.1 — Homework KEY
\documentclass[10pt]{article}
\usepackage{apstats-article}
\usepackage{apstats-key}

%=============================================================================
\begin{document}

\pageheader{Unit 1, Lesson 1.1}{Homework — Answer Key}
\begin{tabularx}{\linewidth}{X X X}
Name:\;\ans{Key} & Date:\; & Period:\;
\end{tabularx}

\vspace{0.1in}

\begin{tcolorbox}[colback=sky, colframe=navy, boxrule=0.8pt, arc=2mm,
    left=3mm, right=3mm, top=1.5mm, bottom=1.5mm]
\small For each news blurb, identify the \textbf{individual}, \textbf{population},
\textbf{sample}, and \textbf{variable(s)}. Write in complete phrases, not single words.
If no sample size is stated, write \textit{not given}.
\end{tcolorbox}

\vspace{0.14in}

% ════════════════════════════════════════════════════════════════════════════
% BLURB 1 KEY
% ════════════════════════════════════════════════════════════════════════════
\begin{blurbbox}[Study 1]{navylight}
\begin{headlinebox}{sky}
\small\textit{``Researchers surveyed 1{,}800 teenagers across 12 countries and
found that 67\% felt social media had a negative effect on their self-image.''}
\end{headlinebox}

\vspace{8pt}
\componenttablekey
  {A teenager (any adolescent in one of the 12 countries studied)}
  {All teenagers (accept: all teenagers in the 12 countries represented)}
  {The 1{,}800 teenagers surveyed}
  {Perception of social media's effect on self-image (categorical: positive / negative / neutral)}

\vspace{6pt}
\begin{teachernote}
\scriptsize
Students may write ``all teenagers in the world'' — accept with a note that the study's
actual scope is the 12 countries sampled. This is a productive ambiguity: the population
the researchers \emph{want} to generalize to may be broader than what their sample
actually represents. Revisit in Unit 3 (sampling methods).
\end{teachernote}
\end{blurbbox}

\vspace{0.12in}

% ════════════════════════════════════════════════════════════════════════════
% BLURB 2 KEY
% ════════════════════════════════════════════════════════════════════════════
\begin{blurbbox}[Study 2]{greenacc}
\begin{headlinebox}{greenbg}
\small\textit{``A study of 250 WNBA players found that athletes who practiced
yoga at least twice a week reported significantly fewer injuries during the season
compared to those who did not.''}
\end{headlinebox}

\vspace{8pt}
\componenttablekey
  {A WNBA player}
  {All WNBA players}
  {The 250 WNBA players in the study}
  {Yoga practice frequency (quantitative: sessions per week); number of injuries during the season (quantitative)}

\vspace{6pt}
\begin{teachernote}
\scriptsize
The WNBA has roughly 140 active roster spots per season — so 250 players likely represents
more than one season or includes players no longer active. Students who notice this are
thinking carefully about feasibility. Worth acknowledging: a near-census of active WNBA
players would actually be achievable, unlike most populations.
\end{teachernote}
\end{blurbbox}

\vspace{0.12in}

% ════════════════════════════════════════════════════════════════════════════
% BLURB 3 KEY
% ════════════════════════════════════════════════════════════════════════════
\begin{blurbbox}[Study 3]{redacc}
\begin{headlinebox}{redbg}
\small\textit{``Scientists collected soil samples from 45 farms in Iowa and
found that fields using cover crops contained 30\% more earthworms per square
meter than fields without cover crops.''}
\end{headlinebox}

\vspace{8pt}
\componenttablekey
  {A farm field (a specific plot of land) --- not a scientist, not a farmer}
  {All farm fields in Iowa (accept: all Midwestern / U.S.\ farm fields)}
  {The 45 farms examined}
  {Number of earthworms per square meter (quantitative); whether cover crops are used (categorical)}

\vspace{6pt}
\begin{teachernote}
\scriptsize
Common errors: writing ``a scientist'' or ``an Iowa farmer'' as the individual.
The individual is whatever each data row describes — here, each row is a field with
a worm count and a cover-crop status. Reinforce: ask ``what did they measure \emph{one of}?''
\end{teachernote}
\end{blurbbox}

\vspace{0.12in}

% ── REFLECTION KEY ────────────────────────────────────────────────────────────
\begin{reflectionbox}
\small Choose \textbf{one} of the three studies above and answer:
\textit{Why might the results of this study not apply to everyone?}
What would make you more --- or less --- confident in the findings?

\vspace{8pt}
\ans{\textbf{Study 1 (model):} The study covered only 12 countries, so results may not
reflect teenagers in countries with different social media cultures or access levels.
Confidence increases if the 12 countries were chosen to represent a broad range of
regions and incomes; it decreases if they were all similar wealthy democracies.}

\vspace{6pt}
\ans{\textbf{Study 2 (model):} Results apply only to elite professional basketball players —
not recreational athletes or players in other sports. Confidence increases if players were
randomly assigned to yoga vs.\ no yoga (an experiment); it decreases if players self-selected
(those who do yoga may already take better care of their bodies).}

\vspace{6pt}
\ans{\textbf{Study 3 (model):} Iowa farms may differ from farms in other states in soil
type, climate, and farming practices, so results may not generalize nationally. Confidence
increases with a larger, more geographically diverse sample.}
\end{reflectionbox}

\vspace{0.12in}

% ── EXTENSION KEY ────────────────────────────────────────────────────────────
\begin{extensionbox}
\small

\textbf{Part A --- U.S.\ Census}

\begin{enumerate}[leftmargin=1.5em, itemsep=6pt]
  \item \textbf{Cost and time:}\\
        \ans{The 2020 U.S.\ Census cost approximately \$14.2 billion and took several years
        to plan, conduct, and process. It employed hundreds of thousands of temporary workers
        to contact households that did not respond online or by mail.}

  \item \textbf{Hard-to-count group:}\\
        \ans{Examples: people experiencing homelessness (no fixed address), undocumented
        immigrants (may fear government contact), people in rural areas (fewer mailings
        reach them), college students (counted at home or at school?). Accept any group
        with a clear explanation.}

  \item \textbf{Why not just use a sample?}\\
        \ans{Census data are used to allocate Congressional seats and federal funding to
        states and counties. A sample would introduce uncertainty into these high-stakes
        decisions. Some applications genuinely require counts of every individual, not estimates.}
\end{enumerate}

\vspace{8pt}
\textbf{Part B --- Preview: Categorical vs.\ Quantitative}

\vspace{6pt}
\renewcommand{\arraystretch}{1.6}
\begin{tabularx}{\linewidth}{@{} >{\bfseries}p{0.08\linewidth}
                                    >{\raggedright\arraybackslash}p{0.44\linewidth}
                                    X @{}}
& \textbf{Variable} & \textbf{Number or category?} \\
\hline
S1 & \ans{Perception of social media's effect on self-image} & \ans{Category} \\
S2 & \ans{Yoga sessions per week; number of injuries}        & \ans{Number (both)} \\
S3 & \ans{Earthworms per sq.\ meter; cover crop use}         & \ans{Number; Category} \\
\end{tabularx}
\end{extensionbox}

\end{document}
